\section{Strong finite energy and finite pivots}
\label{sec:finite-energy}

For $C\subset\Z^d$, let $\cF_C$ denote the sigma-field generated by
$\{X_v:v\in C\}$.  We use the following terminology only for the explicit
condition printed in the definition.

\begin{definition}\label{def:sfe}
A law $\mu$ on $A^{\Z^d}$ has \emph{uniform strong finite energy} with
constant $\eta>0$ if
\[
  \Prob\bigl(X_v=c\mid\cF_{\Z^d\setminus\{v\}}\bigr)\geq\eta
\]
almost surely for every $v\in\Z^d$ and $c\in A$.
\end{definition}

Summing the conditional probabilities over $c$ shows that necessarily
$\eta\leq |A|^{-1}$.  The strength of \cref{def:sfe} is that each coordinate
is conditioned on all the others.  This makes the following iteration
possible.

\begin{lemma}[Conditional block bound]\label{lem:block}
Suppose that $\mu$ has uniform strong finite energy with constant $\eta$.
For every finite $B\subset\Z^d$ and $b\in A^B$,
\begin{equation}\label{eq:block-bound}
  \Prob(X_B=b\mid\cF_{B^c})\geq\eta^{|B|}
  \qquad\text{almost surely}.
\end{equation}
More generally, if $g\colon A^{\Z^d}\to A^B$ is
$\cF_{B^c}$-measurable, then
\begin{equation}\label{eq:random-target}
  \Prob\bigl(X_B=g(X)\mid\cF_{B^c}\bigr)\geq\eta^{|B|}
  \qquad\text{almost surely}.
\end{equation}
\end{lemma}

\begin{proof}
The assertion is immediate for $B=\varnothing$.  Otherwise enumerate
$B=\{v_1,\ldots,v_m\}$ and put
$I_i=\one_{\{X_{v_i}=b_{v_i}\}}$.  Let $H$ be any bounded nonnegative
$\cF_{B^c}$-measurable random variable.  The product
$H I_1\cdots I_{m-1}$ is measurable with respect to
$\cF_{\Z^d\setminus\{v_m\}}$.  Hence
\begin{align*}
 \EE[H I_1\cdots I_m]
 &=\EE\!\left[
   H I_1\cdots I_{m-1}
   \EE(I_m\mid\cF_{\Z^d\setminus\{v_m\}})
 \right]\\
 &\geq \eta\,\EE[H I_1\cdots I_{m-1}].
\end{align*}
Repeating the argument for $v_{m-1},\ldots,v_1$ gives
\[
  \EE[H\one_{\{X_B=b\}}]\geq\eta^m\EE[H].
\]
Testing against all such $H$ is equivalent to \eqref{eq:block-bound}.  In
particular, the argument does not condition on an individual outside
configuration, which could have probability zero.

For the second assertion, $A^B$ is finite and $g$ is
$\cF_{B^c}$-measurable.  Therefore
\begin{align*}
 &\Prob\bigl(X_B=g(X)\mid\cF_{B^c}\bigr)\\
 &\quad=\sum_{b\in A^B}
   \one_{\{g(X)=b\}}\Prob(X_B=b\mid\cF_{B^c})
 \geq \eta^{|B|}.
\end{align*}
\end{proof}

The next proposition packages the only probability step used in both
dynamical arguments.

\begin{proposition}[Measurable site pivots]\label{prop:pivot}
Let $S,T\subset\Z^d$ be finite, let
$\Phi\colon A^T\to A^S$, and let $B\subset T$ satisfy $|B|=|S|$.  Suppose
that, for every fixed pattern on $T\setminus B$, the map
\[
  A^B\longrightarrow A^S,\qquad
  x_B\longmapsto\Phi(x_B,x_{T\setminus B}),
\]
is a bijection.  If $\mu$ has uniform strong finite energy with constant
$\eta$, then for every $a\in A^S$,
\begin{equation}\label{eq:pivot-conditional}
  \Prob\bigl(\Phi(X_T)=a\mid\cF_{B^c}\bigr)
  \geq\eta^{|S|}
  \qquad\text{almost surely}.
\end{equation}
\end{proposition}

\begin{proof}
For each pattern on $T\setminus B$, bijectivity gives a unique pivot pattern
that produces $a$.  These finitely many values define a function
$g_a\colon A^{T\setminus B}\to A^B$.  Thus
\[
   \{\Phi(X_T)=a\}=\{X_B=g_a(X_{T\setminus B})\}.
\]
The target on the right is $\cF_{B^c}$-measurable, so
\eqref{eq:pivot-conditional} follows from
\cref{lem:block}.
\end{proof}

\begin{remark}\label{rem:pivot-normalization}
The exponent in \eqref{eq:pivot-conditional} counts whole input sites.  A
linear-algebra argument that produces several scalar-coordinate pivots on
different sites for one vector-valued output does not meet the hypothesis of
\cref{prop:pivot} with $|B|=|S|$.  This normalization is the obstruction in
\cref{prop:vector-firewall}.
\end{remark}

